\sect{Queries}

The efficient retrieval of information stored in probability distributions has to exploit the available decomposition schemes.
To avoid the instantiation of a distribution based on its decomposition, we directly define deductive reasoning schemes by contractions, which can be executed using the available decomposition.

\subsect{Querying by Functions}

We now formalize queries by retrieving expectations of functions given a distribution specified by probability tensors.
We exploit basis calculus in defining categorical variables $\headvariableof{\exfunction}$ to tensors $\exfunction$, which are enumerating the set $\imageof{\exfunction}$.
More details on this scheme are provided in \charef{cha:basisCalculus}, see \defref{def:functionRelationEncoding} therein.

\begin{definition}
    \label{def:queries}
    The marginal query of a probability distribution $\probat{\shortcatvariables}$ by a query statistic
    \begin{align*}
        \exfunction \defcols \facstates \rightarrow \arbsetof{\exfunction}
    \end{align*}
    is the vector $\probat{\headvariableof{\exfunction}}$, where $\headvariableof{\exfunction}$ is an image enumerating variable (see \charef{cha:basisCalculus} for more details), defined as the contraction
    \begin{align*}
        \probat{\headvariableof{\exfunction}}
        = \contractionof{\probat{\shortcatvariables},\bencodingofat{\exfunction}{\headvariableof{\exfunction},\shortcatvariables}}{\headvariableof{\exfunction}} \, .
    \end{align*}
    If $\arbsetof{\exfunction}\subset\rr$, and the statistic $\exfunction$ is therefore a tensor, we further define the expectation query of $\probwith$ by $\exfunction$ as
    \begin{align*}
        \expectationof{\exfunction} = \contraction{\exfunctionat{\shortcatvariables},\probwith} \, .
    \end{align*}
    Given another query statistic $\secexfunction: \facstates \rightarrow \arbsetof{\secexfunction}$, which image is enumerated by a variable $\headvariableof{\secexfunction}$, the conditional query of the probability distribution $\probat{\shortcatvariables}$ by the statistic $\exfunction$ conditioned on $\secexfunction$ is the matrix $\condprobof{\headvariableof{\exfunction}}{\headvariableof{\secexfunction}}\in\rr^{\cardof{\imageof{\exfunction}}}\otimes \rr^{\cardof{\imageof{\secexfunction}}}$ defined as the normalization
    \begin{align*}
        \condprobof{\headvariableof{\exfunction}}{\headvariableof{\secexfunction}}
        = \normalizationofwrt{
            \probat{\shortcatvariables},\bencodingofat{\exfunction}{\headvariableof{\exfunction},\shortcatvariables},\bencodingofat{\secexfunction}{\headvariableof{\secexfunction},\shortcatvariables}
        }{
            \headvariableof{\exfunction}}{\headvariableof{\secexfunction}
        } \, .
    \end{align*}
\end{definition}

We notice that marginal and conditional queries generalize the schemes of marginalization and conditioning on the distributed variables.
As such, marginal distributions are marginal queries with respect to a identity query statistic acting on the respective variables.
Conditional distributions can be similarly retrieved by two identity statistics, representing the incoming and outgoing variables.
%The query framework of \defref{def:queries} can be used beyond marginal and conditional quereis for the customized information retrieval.

%% Relation of queries and expectation queries
Expectation queries return the expectation of a real-valued feature $\exfunction : \facstates\rightarrow\rr$.
When understanding $\exfunction$ as a random variable given the probability measure $\probat{\shortcatvariables}$, the expectation query returns the expectation of that random variable.
Expectation queries are further contractions of marginal queries with the identity function restricted on the image of $\exfunction$, since
\begin{align*}
    \expectationof{\exfunction}
    = \contraction{\probat{\headvariableof{\exfunction}},\idrestrictedto{\imageof{\exfunction}}{\headvariableof{\exfunction}} } \, .
\end{align*}
This contraction equation follows from the more general \corref{cor:rhoToNormal}, which will be shown in \charef{cha:basisCalculus}.

% Mean parameters
Expectation queries compute mean parameters to a distribution, see \defref{def:meanPolytope}.
Given a statistic $\sstat$ and a distribution $\probwith$, each feature $\sstatcoordinateof{\selindex}$ for $\selindexin$ corresponds with a mean parameter
\begin{align*}
    \expectationof{\exfunction} = \contraction{\probwith,\sstatcoordinateof{\selindex}} = \meanparamat{\indexedselvariable} \, .
\end{align*}

%% Conditional Probabilities and conditional queries
%Conditional probabilities are queries, where the tensors $\exfunction$ and $\secexfunction$ are identity mappings in the respective variable state spaces.
%Conversely, we can understand the conditional query $\condprobof{\headvariableof{\exfunction}}{\headvariableof{\secexfunction}}$ as the conditional probability of a query statistic $\exfunction$ conditioned on $\secexfunction$, of the underlying Markov Network with cores $\{\probtensor, \bencodingof{\exfunction}, \bencodingof{\secexfunction} \}$ and variables $\catvariableof{\exfunction},\catvariableof{\secexfunction}$ besides the variables distributed by $\probtensor$.

%%% Expectations as event queries -> Consistency with $\probat{X=i}$?
%We further denote event queries by
%	\[  \expectationof{\exfunction=\exfunctionimageelement} = \contraction{\probtensor,\bencodingof{\exfunction},\onehotmapof{\exfunctionimageelement}}\]
%where by $\onehotmapof{z}$ be denote the one hot encoding of the state $z$ with respect to some enumeration.
%Let us note that they are further contraction of the queries in \defref{def:queries} since by \theref{the:splittingContractions}
%\begin{align*}
%	 \expectationof{\exfunction=z}
%	& =  \contraction{ \contractionof{\probtensor,\bencodingof{\exfunction}}{\catvariableof{\exfunction}} ,\onehotmapof{z}}\\
%	& =  \contraction{ \probat{\exfunction} ,\onehotmapof{z}} \, .
%\end{align*}

\subsect{Mode Queries}\label{sec:modeQueries}

A different kind of queries are mode queries, which we formalize by the searches of the indices to maximal coordinates in a tensor.

\begin{definition}
    Given a tensor $\hypercorewith$ the mode query is the problem
    \begin{align}
        \argmax_{\shortcatindicesin} \hypercoreat{\indexedshortcatvariables} \, .
    \end{align}
\end{definition}

%MAP relation - Needed?
%Mode queries are often called MAP queries.
%The name MAP is an abbreviation of maximum a-posteriori, since it searches for the most probable state in probability distributions.
%The term a-posteriori originates from a focus on conditional probabilities, which slices with respect to incoming variables represent the a-posteriori distribution given evidence on the incoming variables.
%We will in this work use this term for generic maximum coordinate searches.

% Convex optimization
We can pose mode queries as convex optimization problems.
To this end we recall, that the set of all probability distributions is a convex hull of the one-hot encoded states, that is
\begin{align*}
    \bmrealprobof{\ones}
    = \convhullof{\onehotmapofat{\shortcatindices}{\shortcatvariables}\wcols\shortcatindicesin} \, .
\end{align*}
With this we have
\begin{align*}
    \max_{\shortcatindices} \hypercoreat{\indexedshortcatvariables}
    &= \max_{\shortcatindices} \contraction{\hypercoreat{\shortcatvariables},\onehotmapofat{\shortcatindices}{\shortcatvariables}} \\
    &= \max_{\meanparamat{\shortcatvariables}\in\bmrealprobof{\ones}} \contraction{\hypercoreat{\shortcatvariables},\meanparamat{\shortcatvariables}} \, .
\end{align*}
We note that the maximization over $\bmrealprobof{\ones}$ is a convex optimization problem and the maxima are taken at
\begin{align*}
    \argmax_{\meanparamat{\shortcatvariables}\in\bmrealprobof{\ones}} \contraction{\hypercoreat{\shortcatvariables},\meanparamat{\shortcatvariables}}
    = \convhullof{\onehotmapofat{\shortcatindices}{\shortcatvariables}\wcols\shortcatindices\in\argmax_{\shortcatindices} \hypercoreat{\indexedshortcatvariables}} \, .
\end{align*}
We will further apply this generic trick to approach mode queries when studying inference problems for exponential families in the following sections.

% Restriction to subsets
Let us note, that we can also approach modified mode queries, where we restrict to $\arbset\subset\facstates$, namely
\begin{align*}
    \argmax_{\shortcatindices\in\arbset} \hypercoreat{\indexedshortcatvariables} \, .
\end{align*}
We can choose the base measure by the subset encoding (see \charef{cha:basisCalculus}) of $\arbset$, namely
\begin{align*}
    \basemeasurewith = \sum_{\shortcatindices\in\arbset} \onehotmapofat{\shortcatindices}{\shortcatvariables} \, ,
\end{align*}
and conclude that
\begin{align*}
    \argmax_{\meanparamat{\shortcatvariables}\in\bmrealprobof{\basemeasure}} \contraction{\hypercoreat{\shortcatvariables},\meanparamat{\shortcatvariables}}
    = \convhullof{\onehotmapofat{\shortcatindices}{\shortcatvariables} \wcols \shortcatindices\in\argmax_{\shortcatindices\in\arbset} \hypercoreat{\indexedshortcatvariables}} \, .
\end{align*}



\subsect{Energy Representations}

For exponential families (see \secref{sec:exponentialFamilies}) we have observed, that often energy tensors have feasible tensor network representations, whereas the corresponding probability distributions can get infeasible.
We therefore investigate here schemes to answer queries based on the energy tensor instead of the distribution.
%Let us now interpret a probability tensor at hand as a member of an exponential family (see \secref{sec:exponentialFamilies}), which is always possible when taking the naive exponential family.

\begin{theorem}
    \label{the:energyContractionQueries} % This is a statement about "full" queries.
    Let $\energytensorwith$ be an energy tensor and $\probwith=\normalizationof{\expof{\energytensorwith}}{\shortcatvariables}$ the corresponding distribution.
%	For any probability distribution $\probtensor$ with $\probtensor= \normalizationof{\expof{\energytensorwith}}{\shortcatvariables}$,
    For disjoint subsets $\nodesa,\nodesb \subset [\catorder]$ with $\nodesa\cup\nodesb=[\catorder]$ and any $\catindexof{\nodesb}$ we have
    \begin{align*}
        \condprobof{\catvariableof{\nodesa}}{\indexedcatvariableof{\nodesb}}
        = \normalizationof{
            \expof{\energytensorat{\catvariableof{\nodesa},\indexedcatvariableof{\nodesb}}}
        }{\catvariableof{\nodesa}} \, .
    \end{align*}
\end{theorem}
\begin{proof}
    To show the theorem, we use a generic simplification property of coordinatewise transforms, which we will show as \lemref{lem:coordinatewisetrafoSliceReduction} in \charef{cha:coordinateCalculus} and get
    \begin{align*}
        \contractionof{\expof{\energytensorat{\catvariableof{\nodesa},\catvariableof{\nodesb}}}}{\catvariableof{\nodesa},\indexedcatvariableof{\nodesb}}
        = \contractionof{\expof{\energytensorat{\catvariableof{\nodesa},\indexedcatvariableof{\nodesb}}}}{\catvariableof{\nodesa}}
    \end{align*}
    Based on this we get
    \begin{align*}
        \condprobof{\catvariableof{\nodesa}}{\indexedcatvariableof{\nodesb}}
        &= \normalizationofwrt{\expof{\energytensorat{\catvariableof{\nodesa},\catvariableof{\nodesb}}}}{\catvariableof{\nodesa}}{\indexedcatvariableof{\nodesb}} \\
        &= \frac{\contractionof{\expof{\energytensorat{\catvariableof{\nodesa},\catvariableof{\nodesb}}}}{\catvariableof{\nodesa},\indexedcatvariableof{\nodesb}}}{
            \contractionof{\expof{\energytensorat{\catvariableof{\nodesa},\catvariableof{\nodesb}}}}{\indexedcatvariableof{\nodesb}}
        } \\
        &= \frac{\contractionof{\expof{\energytensorat{\catvariableof{\nodesa},\indexedcatvariableof{\nodesb}}}}{\catvariableof{\nodesa}}}{
            \contraction{\expof{\energytensorat{\catvariableof{\nodesa},\indexedcatvariableof{\nodesb}}}}
        } \\
        &= \normalizationof{
            \expof{\energytensorat{\catvariableof{\nodesa},\indexedcatvariableof{\nodesb}}}
        }{\catvariableof{\nodesa}} \, . \qedhere
    \end{align*}
\end{proof}

% Situation of marginalized variables
Importantly, \theref{the:energyContractionQueries} does not generalize to situations, where $\nodesa\cup\nodesb\neq[\catorder]$, since summation over the indices of the variables $[\catorder]/\nodesa\cup\nodesb$ and contraction do not commute.
In this more generic situation, we would need to sum over exponentiated coordinates, that is
\begin{align*}
    \condprobof{\catvariableof{\nodesa}}{\indexedcatvariableof{\nodesb}}
    = \normalizationof{
        \sum_{\catindexofin{[\catorder]/\nodesa\cup\nodesb}}
        \expof{\energytensorat{\catvariableof{\nodesa},\indexedcatvariableof{\nodesb},\indexedcatvariableof{[\catorder]/\nodesa\cup\nodesb}}}
    }{\catvariableof{\nodesa}} \, .
\end{align*}


% Mode queries
Mode queries on probability distributions in an energy representation can always be reduces to mode queries on the energy tensor.
This is due to the monotonicity of the exponential function, which implies
\begin{align*}
    \argmax_{\shortcatindicesin} \probat{\indexedshortcatvariables}
    &=\argmax_{\shortcatindicesin} \normalizationof{\expof{\energytensorwith}}{\indexedshortcatvariables} \\
    &=\argmax_{\shortcatindicesin} \expof{\energytensorat{\indexedshortcatvariables}} \\
    &=\argmax_{\shortcatindicesin} \energytensorat{\indexedshortcatvariables} \, .
\end{align*}
Since we are only interested in identifying the index of the maximum coordinate, and not its value, we have further dropped the normalization term by partition functions.
When one instead need to get the value of the maximal, the partition function cannot be ignored.